\chapter*{Loch and Waterways}\stepcounter{chapter}\phantomsection\addcontentsline{toc}{chapter}{Loch and Waterways}
\section*{Eas na h-Ùruisg}\phantomsection\addcontentsline{toc}{section}{Eas na h-Ùruisg}
\subsection*{Encounter Type}
{\entryfont Folkloric Investigation and Social Encounter\\
\textbf{Recommended Level:} any\\
\textbf{Difficulty:} Easy}
\subsection*{Setup}
{\entryfont This Encounter can be integrated in your party's adventure in two different ways. Choose whatever suits best during your travelling scene:}
\subsubsection*{The Ùruisge Cascade}
{\entryfont The party comes across a scenic view of a beautiful waterfall. At its base the players can find several plates bearing bannocks in various stages of decay, vessels of sour milk, and wilted flowers. The older the ofering the worse its condition.

Players can find a hidden recess behind the falling water that leads into a small cave. A \textbf{DC 13 Wisdom (Survival) Check} reveals that the cave was once a retreat for a small, humanoid creature - fading foot prints and a disturbed moss bed hints, however, that it has been abandoned for several weeks.

If searched through further the players can find a silver brooch worth 75 GP - but taking it has consequences on the final outcome of this encounter.}
\subsubsection*{The Farms}
{\entryfont Further downstream from the waterfall are standing two farms on either side of the now larger stream. A once great bridge connecting both estates has been washed away - a \textbf{DC 17 Intelligence (Investigation) Checks} reveals that it has been demolished willingly by axe strikes.

\paragraph*{Prosperous Farm} One farmhouse stands within great fields of lush crops - grain bends beneath heavy heads, fruit crowds every branch, and vegetables grow almost grotesquely large. Its mill next to the river works splendidly - seemingly without a break.

Its patriarch greets some strangers with immediate suspicion but to others - that he seems to be probable business partners - he is very welcoming, explaining his superior methods and work to be the basis of his success.

\paragraph*{A somber farm} Very different, this farm stands amid yellowed grain, bare vegetable plots, and dying fruit trees. Its occupants, however, welcome by-travellers and readily admit their fear that they may lose everything.}
\subsection*{Mechanics}
\subsubsection*{The Missing Ùruisge}
{\entryfont The family of the somber farm explain that they have long maintained a friendly relationship with the Ùruisge of the river, leaving it bannock, milk, and occasional flowers. The crops in the region had always prospered in return. However, he creature stopped collecting the offerings several weeks ago, and their harvest began failing soon afterward.}
\subsubsection*{The Prosperous Farm}
{\entryfont The patriarch claims superior seed, irrigation, and hard work explain his harvest. A successful \textbf{DC 14 Wisdom (Insight) Check} reveals that his hostility becomes particularly pronounced whenever the river or neighbouring farm is mentioned.

A \textbf{DC 13 Intelligence (Investigation) Check} around the farm reveals that several irrigation channels converge upon a locked outbuilding near the river.

Beneath it lies a crude stone chamber through which diverted river water flows.

The Ùruisge is imprisoned there.

Iron restraints bind the thin, exhausted creature beside the channel. Empty bowls and signs of prolonged mistreatment surround it. The farmers have discovered that distressing the creature causes sudden surges of unnatural growth in land irrigated by the passing water.

They are not creating fertility.

They are drawing it away from the rest of the valley.}
\subsection*{Possible Solutions}
\subsubsection*{Free the Ùruisge}
{\entryfont Breaking or unlocking its restraints requires a DC 15 Strength Check or a DC 14 Dexterity Check using thieves' tools.

The terrified creature initially assumes its rescuers intend to exploit it as well. Calming it requires a \textbf{DC 14 Wisdom (Animal Handling)}, \textbf{Charisma (Persuasion)}, or \textbf{Charisma (Performance) Check}. Offering fresh bannock or milk grants advantage.

Once reassured, it immediately seeks the river and returns toward the waterfall.}
\subsubsection*{Confront the Farmers}
{\entryfont Evidence of the captive creature makes the family's actions difficult to conceal. The patriarch may surrender if confronted by overwhelming evidence or threatened with exposure.

A successful \textbf{DC 15 Charisma (Intimidation or Persuasion) Check} convinces the remaining family members to stand down and release the Ùruisge themselves.}
\subsubsection*{Force}
{\entryfont If attacked or caught freeing the creature, the patriarch and several farmhands may fight.

They defend the farm rather than pursue fleeing characters and surrender once defeat is inevitable.}
\subsection*{Possible Complications}
\subsubsection*{The Stolen Brooch}
{\entryfont If the characters took the silver brooch from the waterfall cave, the Ùruisge notices its absence upon returning home.

If it is returned voluntarily, there is no penalty.

If it is kept, the creature still restores the river but refuses any personal blessing or reward.}
\subsubsection*{Harming the Ùruisge}
{\entryfont Any character who deliberately harms the captive Ùruisge has disadvantage on checks made to calm it. If seriously injured after being freed, it flees the valley entirely.}
\subsubsection*{Killing the Captors}
{\entryfont The Ùruisge does not mourn its tormentors, but unnecessary cruelty makes it wary of its rescuers. Characters who slaughter surrendering farmers cannot receive its blessing.}
\subsection*{Outcomes}
\subsubsection*{The River Restored}
{\entryfont If freed and treated respectfully, the Ùruisge returns to the waterfall. Over the following days the unnatural abundance of the second farm fades, while water, colour, and life gradually return to the surrounding fields.

The older family resumes leaving offerings - this time with considerably less concern for keeping the reason secret.}
\subsubsection*{Rewards for Proper Conduct}
{\entryfont If the party frees the Ùruisge without stealing from it or committing needless cruelty, it may grant one or more of the following:
\begin{itemize}
	\item \textbf{Blessing of the River:} Each character gains advantage on their next saving throw against exhaustion or an environmental hazard involving water.
	\item \textbf{River's Bounty:} The grateful farming family provides food and lodging whenever the characters pass through the valley.
	\item \textbf{Hidden Gift:} The Ùruisge retrieves a small object from behind the falls: a pearl, polished river-stone charm, or similar treasure worth 75-100 GP.
\end{itemize}
If the silver brooch was stolen and kept, the valley is restored, but the Ùruisge disappears behind the waterfall without offering a reward.}

\clearpage

\section*{Kelpie Bridge}\phantomsection\addcontentsline{toc}{section}{Kelpie Bridge}
\subsection*{Encounter Type}
{\entryfont Comedic Fey Social Encounter with optional Escalating Fey Retribution\\
\textbf{Recommended Level:} 4 - 6\\
\textbf{Difficulty:} Moderate (Deadly if players are isolated)}
\subsection*{Setup}
{\entryfont As the party approaches a modest wooden bridge spanning dark water, they find it occupied by two unusually small black horses.

They are clearly young — slightly too gangly, slightly too small, their manes overly dramatic for their size. They stand in the exact center of the bridge, trying very hard to look imposing.}
\begin{DndReadAloud}
	Two small black horses stand squarely in the middle of the bridge.

	They are... not large. Not majestic. In fact, they look barely old enough to be away from their mothers. But they stand rigidly, chins lifted, eyes glowing faintly red as if this is Very Serious Business.

	One stamps a hoof.

	"Halt!" it declares.

	The other quickly adds, "Yes! Halt! None may cross without answering our question!"

	There is a pause.

	The first whispers, loudly, "Which one was it?"
\end{DndReadAloud}

\subsection*{Mechanics}
{\entryfont The young kelpies attempt to present a riddle, but they cannot agree which one it should be. They cycle through options:
\begin{itemize}
	\item "What weighs more — a secret kept or a promise broken?"
	\item "If you drown someone politely, is it still rude?"
	\item "Is fog just shy rain?"
	\item "How many dramatic pauses are too many?"
\end{itemize}
They argue mid-sentence. They interrupt each other. They occasionally forget what they were asking.

There is no \textbf{correct answer}.

Any thoughtful, absurd, poetic, or dramatic response impresses them equally — provided it is delivered with flair.}

\subsection*{Possible Solutions}
{\entryfont If the party engages playfully, the encounter becomes entirely social.

The kelpies adore:
\begin{itemize}
	\item Embarrassing personal stories (the more dramatic, the better).
	\item Songs, instruments, theatrical performances.
	\item Excessive flattery ("Your mane catches the light like liquid night.").
	\item Polished metal mirrors (they are vain and fascinated by their reflections).
	\item Being told they are terrifying and important guardians.
	\item A successful interaction requires:
\end{itemize}
2 - 3 successful skill checks (DC 14), such as:
\begin{itemize}
	\item Performance
	\item Persuasion
	\item Deception (if amusing)
	\item Animal Handling (if gentle)
	\item Arcana (to treat them with proper fey etiquette)
\end{itemize}
On success, the young kelpies become delighted.}

\subsection*{Possible Complications}
{\entryfont If the party attacks the young kelpies the tone shifts instantly. The young kelpies shriek and vanish beneath the water. The bridge shudders.}

\subsubsection*{Environmental Escalation}
{\entryfont The planks become slick and writhe like something alive (DC 14 Dex save or fall prone). Water lashes upward in snapping arcs. Hoofbeats echo from nowhere. A spectral horse made of river mist charges across the bridge.

One character (random or closest to edge) must succeed on a DC 15 Strength save or be dragged into the river by the spectral steed.

From the dark water rise two full-grown \hyperref[monster:Kelpie]{\LinkFont{Kelpies}} — larger, gaunt, eyes burning with predatory intelligence.

The adults fight intelligently:
\begin{itemize}
	\item Attempting to drag one or two characters into the water.
	\item Targeting isolated or prone characters.
	\item Retreating into water between strikes.
\end{itemize}
The young kelpies remain hidden below unless the adults are slain.}

\subsection*{Outcomes}
\subsubsection*{Rewards for Proper Conduct}
{\entryfont The young kelpies may grant one or more of the following:
\begin{itemize}
	\item \textbf{Safe Passage} The water stills and the bridge becomes firm and dry.
	\item \textbf{Fey Favour} One character gains advantage on their next Persuasion check involving fey creatures.
	\item \textbf{Hidden Treasure} They proudly reveal items snagged beneath the bridge:
	\begin{itemize}
		\item 30 - 60 GP in trinkets
		\item A silver brooch
		\item A minor potion
		\item An oddly intact waterproof scroll tube
	\end{itemize}
\end{itemize}
They insist it is a "tribute", not "stuff we found".}